\SourceLeaf{450}%
\par\Needspace{9\baselineskip}\begin{center}\textbf{BEMERKUNGEN.}\end{center}
\noindent\hspace*{1.2em}Die S.\ 397 dieses Bandes wiedergegebenen Mittheilungen an \textsc{Olbers} vom 9.\ October 1825 zeigen, dass \textsc{Gauss} zu dieser Zeit begonnen hat, die Untersuchungen zu Papier zu bringen, die den Gegenstand seiner \textit{Disquisitiones generales circa superficies curvas} vom Jahre 1827 bilden. Das vorstehend abgedruckte Fragment: \textit{Neue allgemeine Untersuchungen über die krummen Flächen} unterscheidet sich von den \textit{Disquisitiones} nicht nur durch den geringeren Umfang des Stoffes, sondern auch durch die Art der Behandlung und die Anordnung der Sätze. Dort nimmt \textsc{Gauss} an, dass die rechtwinkligen Coordinaten $x,y,z$ eines Punktes der Fläche als Functionen von irgend zwei unabhängigen Veränderlichen $p$ und $q$ dargestellt seien, während er hier als neue Veränderliche die geodätischen Coordinaten $s$ und $\theta$ wählt. Hier beweist er zuerst den Satz, dass die Summe der drei Winkel eines Dreiecks, das auf einer beliebigen krummen Fläche durch kürzeste Linien gebildet wird, sich von $180^{0}$ um den Inhalt des Dreiecks unterscheidet, das ihm bei der Abbildung durch parallele Normalen auf der Hülfskugel vom Radius $1$ entspricht, und leitet daraus durch einfache geometrische Betrachtungen das fundamentale Theorem her, dass »bei der Übertragung der Flächen durch Abwickelung das Krümmungsmass an jeder Stelle unverändert bleibt«, während er dort zuerst in \S\ 12 zeigt, dass das Krümmungsmass sich allein durch die drei Grössen $E,F,G$ und deren Ableitungen nach $p$ und $q$ ausdrücken lässt, woraus der Satz über die Erhaltung des Krümmungsmasses als Corollar folgt, und erst viel später, in \S\ 20, ganz unabhängig davon den Satz über die Winkelsumme eines geodätischen Dreiecks beweist.
\begin{flushright}\textsc{Stäckel}.\end{flushright}
